\documentclass[11pt,a4paper]{article}
\usepackage[T1]{fontenc}
\usepackage[utf8]{inputenc}
\usepackage{lmodern,amsmath,amssymb,amsthm,booktabs,array}
\usepackage[margin=25mm]{geometry}
\usepackage{microtype}
\usepackage{xurl}
\usepackage[colorlinks=true,linkcolor=black,citecolor=blue,urlcolor=blue]{hyperref}
\hypersetup{
 pdftitle={An Exact Extension and an Infinite Family for OEIS A249106},
 pdfauthor={Carlo Corti},
 pdfsubject={Exact arithmetic-derivative search and an infinite family},
 pdfkeywords={arithmetic derivative, Pythagorean sextuple, OEIS A249106}
}
\newtheorem{theorem}{Theorem}
\newtheorem{lemma}[theorem]{Lemma}
\newtheorem{proposition}[theorem]{Proposition}
\newtheorem{corollary}[theorem]{Corollary}
\theoremstyle{remark}
\newtheorem{remark}[theorem]{Remark}
\newcommand{\D}{\operatorname{D}}
\newcommand{\N}{\mathbb{N}}
\title{An Exact Extension and an Infinite Family for OEIS A249106}
\author{Carlo Corti\\Independent researcher\\\texttt{carlo.corti@outlook.com}}
\date{October 2026}
\begin{document}
\maketitle
\begin{abstract}
OEIS A249106 lists integers whose squares, together with the squares of their
first four arithmetic derivatives, sum to a square. An exact C enumeration of
the inclusive interval $[434325392,10^{10}]$, comprising $9565674609$ integers,
found precisely $1102826289$ and $3945594513$. Relative to the accepted
five-term OEIS prefix, these are the sixth and seventh terms. Their positive
witnesses have independent trial-division checks; the exclusion of all other
inputs relies on the reviewed enumeration algorithm and its complete segment
records, rather than a second independent scan. Separately, the fixed-point
identity $\D(p^p)=p^p$ yields the family $4791p^p$ for every prime $p$.
An explicit seed calculation and an induction prove membership and hence
infinitude, without extrapolating the finite search. A general seed
criterion also explains analogous constructions at other chain lengths,
including $128p^p$ for A230543. The fixed-point machinery is historical;
absolute priority for the seed specializations is not asserted.
Combining the A249106 family
with the enumeration bounds the eighth term by
$10^{10}<a(8)\le1366928213897301$; its exact value remains undetermined.
\end{abstract}
\noindent\textbf{MSC 2020:} 11A25, 11D09, 11Y55.\\
\textbf{Keywords:} arithmetic derivative; Pythagorean sextuple; exact enumeration;
segmented sieve; OEIS A249106.

\section{Definition and context}
The arithmetic derivative on nonnegative integers is determined by
$\D(0)=\D(1)=0$, $\D(p)=1$ for primes $p$, and the product rule.
For $x\ge2$, unique factorization gives
\begin{equation}\label{eq:derivative}
 x=\prod_p p^{e_p},\qquad \D(x)=\sum_{p\mid x}e_p\frac{x}{p}.
\end{equation}
These elementary properties and the fixed points $p^p$ are discussed by
Ufnarovski and \AA hlander~\cite[Theorem 6]{ufnarovski} and in OEIS A003415~\cite{derivative}.
The map is not additive. Iterated derivatives below always mean composition.
Put $x_j=\D^j(n)$ for $0\le j\le4$, with $x_0=n$, and define
\begin{equation}\label{eq:predicate}
 Q(n)=\sum_{j=0}^4 x_j^2.
\end{equation}
A249106 is the increasing enumeration, with index offset one, of the integers
$n\ge2$ for which $Q(n)=r^2$ for an integer $r\ge0$~\cite{sequence}.
The restriction $n\ge2$ follows the entry's enumeration convention:
$n=1$ would satisfy the bare equation trivially, since $Q(1)=1$.
Its Pythagorean sextuple is $(x_0,x_1,x_2,x_3,x_4,r)$: five chain coordinates
and the root, with no primitivity requirement.

Paolo P. Lava introduced the sequence on October 21, 2014. The entry snapshot
accepted for this project credits Ray Chandler with the fourth term on
December 23, 2016 and the fifth on January 11, 2017. That snapshot has the
five-term prefix
\begin{equation}\label{eq:baseline}
19164,\ 129357,\ 14971875,\ 45316123,\ 434325391.
\end{equation}
The archived entry snapshot agrees with the public OEIS entry, whose
published prefix is~\eqref{eq:baseline}. The two additional terms reported
here have not been submitted to OEIS.
The prefix was checked at its five positive inputs; a full enumeration of
$[2,434325391]$ was not repeated.

Related entries concern different chain lengths or primitivity conditions:
A210503 imposes a primitive three-coordinate condition; A230543 concerns
four-coordinate tuples without that requirement. A249105, A249107 and A249110
concern five-, seven- and ten-coordinate
Pythagorean tuples~\cite{related}. They provide context rather than inclusion
relations with A249106. General polynomial parametrizations of Pythagorean
sextuples~\cite{frisch} do not enforce the successive arithmetic-derivative
constraints in~\eqref{eq:predicate}.

The results here have two distinct sources. The finite search extends the
indexed prefix through seven and supplies an exclusion bound. Subsequent
analysis of four scaled witnesses gives an infinite family by proof. The
fixed-point identity used in that proof is prior arithmetic-derivative theory;
no absolute priority claim is made for its specialization to this seed.
The project artifacts are publicly available from the GitHub repository and
Zenodo archive cited in reference~\cite{project}. The designated Zenodo DOI
identifies the release series rather than one immutable deposited version.

\section{Finite enumeration and witnesses}
The search began immediately after the accepted fifth term and stopped at
$10^{10}$. Seven completed segments partition the interval, as shown in
Table~\ref{tab:coverage}. The two phases used blocks of $2^{20}$ inputs
and eight threads. Phase 1 consisted of one segment configured for at most
540 blocks and clipped at $10^9$. Phase 2 used segments of 1500 blocks,
with the last clipped at $10^{10}$.
\begin{table}[htbp]
\centering
\begin{tabular}{crrr}
\toprule
Phase/segment & First input & Last input & Hits\\
\midrule
1/0 & 434325392 & 1000000000 & 0\\
2/0 & 1000000001 & 2572864000 & 1\\
2/1 & 2572864001 & 4145728000 & 1\\
2/2 & 4145728001 & 5718592000 & 0\\
2/3 & 5718592001 & 7291456000 & 0\\
2/4 & 7291456001 & 8864320000 & 0\\
2/5 & 8864320001 & 10000000000 & 0\\
\bottomrule
\end{tabular}
\caption{Completed inclusive segments of the finite search. The hit counts
come from the enumeration records; the segment checks are not independent
re-enumerations of non-hits.}\label{tab:coverage}
\end{table}

The interval cardinality is exactly
\begin{equation}\label{eq:count}
10^{10}-434325392+1=9565674609.
\end{equation}
Table~\ref{tab:witnesses} gives the entire consolidated prefix and all witness
coordinates. Its first five rows derive from previously validated outputs of
an earlier C build; the mathematical source regions are byte-identical in
the final campaign build. The last two rows are campaign outputs that also
agree with independently computed Python witnesses.
\begin{table}[htbp]
\centering\scriptsize
\setlength{\tabcolsep}{3pt}
\begin{tabular}{rrrrrrr}
\toprule
$i$ & $a(i)$ & $\D(a(i))$ & $\D^2(a(i))$ & $\D^3(a(i))$ & $\D^4(a(i))$ & $r$\\
\midrule
1 & 19164 & 25564 & 31848 & 58412 & 61916 & 96336\\
2 & 129357 & 172557 & 214974 & 394281 & 417933 & 650268\\
3 & 14971875 & 19971875 & 24881250 & 45634375 & 48371875 & 75262500\\
4 & 45316123 & 198116 & 198120 & 419644 & 419648 & 45320875\\
5 & 434325391 & 3315592 & 5447100 & 9543040 & 36609088 & 436016633\\
6 & 1102826289 & 411955468 & 411955472 & 830107536 & 1973010320 & 2477395545\\
7 & 3945594513 & 5263263313 & 6557049366 & 12026198429 & 12747622097 & 19834209612\\
\bottomrule
\end{tabular}
\caption{Seven indexed witnesses. In every row the sum of squares of the
five chain coordinates equals $r^2$.}\label{tab:witnesses}
\end{table}

\begin{theorem}[Finite computational result]\label{thm:finite}
Subject to the accepted prefix~\eqref{eq:baseline} and the correctness and
complete execution of the enumeration described in Section~\ref{sec:algorithm},
the sixth and seventh terms are
\[
a(6)=1102826289,\qquad a(7)=3945594513,
\]
and there is no further member at most $10^{10}$.
\end{theorem}
\begin{proof}[Computational evidence]
The seven completed records have adjacent endpoints and total cardinality
\eqref{eq:count}. The program audits adjacency within each phase separately;
the join from $10^9$ to $10^9+1$ is checked externally against the two
configurations and the completed records. Their combined hit list consists exactly of the two values
displayed. Both positive chains were separately recomputed by exhaustive
trial division and integer square-root tests. All remaining inputs are
excluded by the C enumeration, whose coverage and arithmetic invariants are
given in Section~\ref{sec:algorithm}. File-integrity and native verification checks establish the
consistency of the completed records, but do not constitute an independent
second enumeration. Combining this finite result with the accepted earlier
prefix gives the asserted indices.
\end{proof}

\section{An unconditional infinite family for A249106}\label{sec:family}
\subsection{From observations to a statement for every prime}
The first three terms and the seventh have the respective multipliers
$2^2,3^3,5^5,7^7$ of $4791=3\cdot1597$. This observation admits an exact
extension to every prime, using the known identity $\D(p^p)=p^p$.
The observations suggested the seed; they do not prove the universal
statement. The proof uses only the defining product rule, four small
factorizations and one integer square identity. It applies to every prime,
including primes beyond any computational search limit.

\subsection{The historical fixed-point mechanism}
For a prime $p$ and integer $k\ge1$, the product rule gives
$\D(p^k)=kp^{k-1}$, hence $\D(p^p)=p^p$.
These fixed points are established prior theory~\cite[Theorem 6]{ufnarovski}.
The transition $T(m)=m+\D(m)$ is itself OEIS A129283~\cite{auxiliary}.
Its connection with multiplication by $p^p$ was recorded in A003415
by R. Zumkeller on April 7, 2007~\cite{derivative}.
Its extension to a whole derivative chain follows by induction.

\begin{lemma}\label{lem:scaling}
Let $t$ be a positive integer with $\D(t)=t$, and put $T(m)=m+\D(m)$
on nonnegative integers. For every $m\ge0$ and $j\ge0$,
\begin{equation}\label{eq:scaling}
 \D^j(tm)=tT^j(m).
\end{equation}
\end{lemma}
\begin{proof}
The case $j=0$ is immediate. If it holds at $j$, the product rule gives
\[
\D\bigl(tT^j(m)\bigr)=\D(t)T^j(m)+t\D\bigl(T^j(m)\bigr)
=tT^{j+1}(m).
\]
This proves the induction. No coprimality condition is needed.
\end{proof}
For $t=p^p$ and $m>1$, the one-step identity appears in
Ufnarovski and \AA hlander~\cite[Theorem 4]{ufnarovski}, whose divergence
statement already uses iteration of that identity. By their Theorem 6,
the positive fixed points $t$ in Lemma~\ref{lem:scaling} are exactly $p^p$.
The lemma states the iterate formula explicitly and includes $m=0,1$.
Since $\D$ is not additive, $T^j$ is a composition,
not a binomial expansion in powers of $\D$.
For example,
$T^2(m)=m+\D(m)+\D(m+\D(m))$; replacing the last derivative by a sum
of derivatives would be unjustified. Shared factors are allowed:
$p=3$ divides the seed, and other primes can divide later seed coordinates.

\subsection{The seed and the unconditional family}
\begin{theorem}[Unconditional family]\label{thm:family}
For every prime $p$, the integer $4791p^p$ belongs to A249106, with witness
\begin{equation}\label{eq:family}
p^p(4791,6391,7962,14603,15479,24084).
\end{equation}
In particular, A249106 is infinite.
\end{theorem}
\begin{proof}
The following four factorizations give the seed calculation directly:
\begin{align*}
4791&=3\cdot1597, &\D(4791)&=1597+3=1600,\\
6391&=7\cdot11\cdot83, &\D(6391)&=913+581+77=1571,\\
7962&=2\cdot3\cdot1327, &\D(7962)&=3981+2654+6=6641,\\
14603&=17\cdot859, &\D(14603)&=859+17=876.
\end{align*}
All displayed factors are prime, as can be checked by trial division through
their integer square roots. Hence successive applications of $T$ yield
$(4791,6391,7962,14603,15479)$, and
\begin{equation}\label{eq:seed}
4791^2+6391^2+7962^2+14603^2+15479^2=580039056=24084^2.
\end{equation}
For prime $p$, equation~\eqref{eq:derivative} gives $\D(p^p)=p^p$.
Lemma~\ref{lem:scaling} and~\eqref{eq:seed} prove~\eqref{eq:family}.
Explicitly, the four derivatives of $4791t$, where $t=p^p$, are
$6391t$, $7962t$, $14603t$ and $15479t$, so
\[
 Q(4791p^p)=p^{2p}\,580039056=(24084p^p)^2.
\]
This is exactly the defining condition of A249106.
The integers $4791p^p$ are distinct and unbounded as $p$ runs over the
infinitely many primes, proving infinitude.
\end{proof}

If $q>p\ge2$, then $q^q>p^q>p^p$, so increasing prime parameters
give increasing family members.
Theorem~\ref{thm:family} has no hypothesis about the accepted prefix,
completion of the C campaign, a primality conjecture, machine-word sizes
or availability of a public software deposit. Its validity follows from
the stated definition of $\D$. The indexed finite result in
Theorem~\ref{thm:finite} has different dependencies.

\subsection{Primitivity, scope and unknown indices}
The seed sextuple is primitive: already $\gcd(4791,6391)=1$. Thus the common
divisor of all six coordinates in~\eqref{eq:family} is exactly $p^p$.
For rows 1 through 7 of Table~\ref{tab:witnesses}, the six-coordinate greatest
common divisors are respectively
\[
4,\ 27,\ 3125,\ 1,\ 1,\ 1,\ 823543.
\]
In particular, rows 4, 5 and 6 are primitive witnesses outside this scaled
family. Their first coordinates are not divisible by $1597$.
Neither that divisibility nor the family is an exhaustive search filter.
The primes $2,3,5,7$ yield indices $1,2,3,7$, so the prime parameter is
not a formula for the sequence index.

The product-rule proof also shows why arbitrary multipliers cannot be
substituted for $p^p$: in general
$\D(sm)=s\D(m)+m\D(s)$, so the seed transition is $T$ only when
$\D(s)=s$. The theorem supplies a sufficient family, not a classification
of every multiplier or every member.
The ratios $\D^j(4791p^p)/(4791p^p)$, $0\le j\le4$, are independent
of $p$, and the root/input ratio is $24084/4791$. The entire witness
changes by a common scale, while its normalized shape stays fixed.

Table~\ref{tab:unindexed} gives four further consequences.
Here $n_p$ is indexed by its prime parameter, not $a(p)$ and not an
assigned OEIS index. Each root is $r_p=24084p^p$ and the chain is
given by~\eqref{eq:family}; the intervening intervals have not been enumerated.
\begin{table}[htbp]
\centering
\begin{tabular}{rr}
\toprule
$p$ & $n_p=4791p^p$\\
\midrule
11 & 1366928213897301\\
13 & 1451074635683484123\\
17 & 3963308094697439437172007\\
19 & 9478608570268562405492983389\\
\bottomrule
\end{tabular}
\caption{Four proven members without assigned sequence indices.
These symbolic consequences do not extend the indexed b-file.}
\label{tab:unindexed}
\end{table}

\subsection{A bound without identification of the next term}
\begin{corollary}\label{cor:eighth}
The upper bound $a(8)\le4791\cdot11^{11}$ holds independently of the
finite exclusion and completeness of the accepted baseline. With the
finite-result assumptions of Theorem~\ref{thm:finite}, the combined bounds are
\begin{equation}\label{eq:eighth}
10^{10}<a(8)\le4791\cdot11^{11}=1366928213897301.
\end{equation}
\end{corollary}
\begin{proof}
The seven distinct positive members in Table~\ref{tab:witnesses}, whose
chains were recomputed by Python trial division as described in
Section~\ref{sec:validation}, and the
$p=11$ member of Theorem~\ref{thm:family} give eight distinct members
at most $4791\cdot11^{11}$. This proves the upper bound without assuming
that the table contains every earlier member. The $p=11$ member exceeds
$10^{10}$ and has root $24084\cdot11^{11}=6871446274995324$.
The lower bound follows from Theorem~\ref{thm:finite}.
\end{proof}
This upper-bound witness is derived symbolically, has no assigned sequence
index, and was not found by extending the computational search. It even
exceeds the implementation's conservative input cap discussed below.

\section{Related families and the scope of the contribution}\label{sec:related}
\subsection{A criterion for any derivative-chain length}
The same proof works for other chain lengths, but each requires its own
seed identity. For integers $h\ge1$ and $m\ge1$, define
\[
 Q_h(n)=\sum_{j=0}^h\bigl(\D^j(n)\bigr)^2,\qquad
 S_h(m)=\sum_{j=0}^h\bigl(T^j(m)\bigr)^2.
\]
Lemma~\ref{lem:scaling} gives, for every prime $p$,
\begin{equation}\label{eq:general-seed}
 Q_h(mp^p)=p^{2p}S_h(m).
\end{equation}
Thus $mp^p$ has a square sum at depth $h$ if and only if $S_h(m)$
is an integer square. For the converse, if $r^2=t^2S_h(m)$ with
$t=p^p$, comparison of prime exponents shows $t\mid r$, so
$S_h(m)=(r/t)^2$. The forward implication follows by multiplying
an integer seed root by $t$. One exact seed calculation therefore proves
membership for every prime; failure rules out all prime parameters for
that seed and depth. This elementary criterion is derived here from the
known product rule, without a claim that the scaling principle is new.
Here $Q=Q_4$, and for A249106 the successful test is $S_4(4791)=24084^2$.

The seed also has a conditional minimality property. If $m\ge1$ and
$S_4(m)$ is a square, the case $p=2$ gives the member $4m$ of A249106.
Thus, assuming the accepted first term $a(1)=19164$ is minimal,
$4m\ge19164$ and $m\ge4791$. Since $4791$ passes the seed test, it is
then the least depth-four seed. Similarly, assuming the recorded first
member $512$ of A230543 is minimal, its least depth-two seed is $128$:
every such seed yields a member $4m\ge512$, and $128$ passes the test
in Section~\ref{sec:quadruple}. These are deductions from the respective
first-term assumptions, not new independent exclusions of smaller inputs.

\subsection{Established analogues in related OEIS entries}
For $m=1$, all iterates of $T$ equal one and $S_h(1)=h+1$.
Whenever $h+1=b^2$ for an integer $b$, every $p^p$ has a constant
derivative chain and root $bp^p$. In particular,
\[
 Q_3(p^p)=(2p^p)^2,\qquad Q_8(p^p)=(3p^p)^2.
\]
The inclusions of A051674, the sequence $p^p$ for prime $p$, are already
recorded in A249105 and A249110~\cite{related,auxiliary}.
They prove infinitude of those sequences without a prime-pair conjecture.
At depth four the seed $m=1$ fails because $S_4(1)=5$ is not a square;
the nonconstant seed $4791$ supplies the needed replacement.

A210503 has a primitive condition. Its entry records the twin-prime
construction~\cite{related}: when $p$ and $p+2$ are primes,
\[
 n=p(p+2),\qquad \D(n)=2(p+1),\qquad
 n^2+\D(n)^2=(n+2)^2.
\]
Here $p$ and $n$ are odd, and neither $p$ nor $p+2$ divides
$2(p+1)$, so $\gcd(n,\D(n))=1$ and the triple is primitive.
The identity follows from $n=(p+1)^2-1$.
Each prime pair gives a member, but infinitude from this construction
would require infinitely many twin primes.
The scaled witnesses of Theorem~\ref{thm:family} have common divisor
$p^p$ and cannot provide primitive triples simply by dropping coordinates.

\subsection{An explicit comparison for A230543}\label{sec:quadruple}
A230543 uses depth two with no primitivity restriction~\cite{related}.
Since $128=2^7$ and $576=2^6\,3^2$,
\[
 \D(128)=448,\quad \D(576)=2112,\quad
 (T^0(128),T(128),T^2(128))=(128,576,2688).
\]
The square identity
\[
 128^2+576^2+2688^2=7573504=2752^2
\]
and~\eqref{eq:general-seed} prove that every prime $p$ gives the
A230543 witness $p^p(128,576,2688,2752)$.
The cases $p=2,3,5,7$ give $512,3456,400000,105413504$, all present
in the consulted entry. This family is derived in the comparison:
the consulted entry does not state the generating formula, and no
priority claim is made for it. Its membership and infinitude follow
from the proof independently of the recorded examples.

For A249107 (depth five), the consulted entry supplies finite examples
but no analogous generating family. This describes that entry, not the
absence of a theorem elsewhere.
Different depths give no inclusion into A249106 merely by deleting or
appending derivatives.

\subsection{Prior art and what is established here}
The historical ingredients are the fixed points $p^p$, the identity
for $\D(mp^p)$ and the transition $m+\D(m)$
~\cite{ufnarovski,derivative,auxiliary}.
Kovi\v{c} also discusses fixed points and higher derivatives in a broader
differential-equation setting~\cite[Sections 1.2 and 2.1]{kovic}.
The contribution made explicit in Theorem~\ref{thm:family} is the
specialization to the verified seed $4791$, its complete chain witness
and the resulting unconditional infinitude of A249106.
It was proved after the finite campaign, rather than used as a filter.

A targeted primary-source and web search on October 3, 2026 did not
identify an earlier explicit statement of the family $4791p^p$ or its
displayed seed witness. This limited search does not establish absolute
priority: indexing gaps and literature unavailable in full text can hide
precedents. In particular, the 2013 Hoepli book by G. Balzarotti and
P. P. Lava, \emph{La derivata aritmetica}, listed in A003415~\cite{derivative},
was not available to this project in full text and was not consulted.
The family is therefore presented with its proof and known
antecedents, without asserting a first discovery of the underlying
mechanism or excluding an earlier occurrence of this particular family.

\section{Exact enumeration algorithm}\label{sec:algorithm}
The campaign used a C17 program with an unsigned 128-bit integer extension,
GMP integer arithmetic, OpenMP, and OpenSSL digest routines. The final code
does not use FLINT. Its mathematical stages are segmented evaluation of the
first derivative, exact factorization of the next three operands, and an
exact square test. Operational segments group consecutive blocks; the
mathematical processing within each block is independent of this grouping.

\subsection{Consecutive inputs and the first derivative}
For each operational segment the program precomputes the primes through the
integer square root of the segment upper endpoint. Each block processes
only the table prefix satisfying $p^2\le H$, where $H$ is the block upper
endpoint. Within a block, a residual array
starts at the original inputs $n$, and a derivative accumulator starts at
zero. For each processed prime $p$, all block multiples of $p$ are visited;
the full multiplicity $e$ is removed from their residuals, and $e(n/p)$ is
added to their accumulators. At every stage, the product of the residual
and the removed prime powers is $n$, and the accumulator is precisely the
contribution of those removed powers in~\eqref{eq:derivative}.

After all primes through the block square root have been processed, any
residual greater than one is prime: a composite residual would have a prime
factor no larger than its square root, hence no larger than that of $n$,
and would already have been divided out. Its remaining contribution $n/q$
is then added. This proves exact evaluation of $\D(n)$ for every block input.

The implemented early exclusion is the prime-input case. For $n\ge2$,
$\D(n)=1$ holds exactly when $n$ is prime. Composite inputs have at least
two prime factors counted with multiplicity and derivative greater than one.
For a prime, the chain is $(n,1,0,0,0)$, and
$n^2<n^2+1<(n+1)^2$. Such inputs cannot qualify. The code uses no filter
requiring the factor $1597$ and no separately implemented congruence filter.

\subsection{Derived operands and deterministic completion}
Each of $x_1,x_2,x_3$ is factored to obtain the next derivative. The word-size
primality routine applies small-prime divisions followed by the strong
probable-prime test with bases
\[
2,\ 325,\ 9375,\ 28178,\ 450775,\ 9780504,\ 1795265022.
\]
These seven bases are sufficient for deterministic primality below $2^{64}$,
as recorded by Fori\v{s}ek and Jan\v{c}ina~\cite[Theorem 3]{forisek}.
Each base is reduced modulo the operand, with zero residues skipped.
Modular products use unsigned 128-bit arithmetic, so they are exact for
word-size operands.

The splitter first checks small divisors through 97. It then tries 32
Pollard-rho-style trajectories~\cite{pollard}, with the recurrence $z\mapsto z^2+c\pmod n$,
Floyd cycle comparison and greatest common divisors, using at most $100000$
steps per trajectory. This is an acceleration rather than a correctness
assumption. If the trajectories do not split a known composite, exhaustive
odd trial division from 101 through its square root supplies a deterministic
fallback. The condition $d\le n/d$ avoids overflow in the loop bound.

A proper factor splits the problem recursively into two smaller positive
operands. Every accepted leaf is checked prime, all multiplicities are
collected, and the reconstructed product is required to equal the original
operand. There are at most 63 prime factors counted with multiplicity for
a positive operand below $2^{64}$, so the 64-entry factor array suffices.
Reconstruction multiplies only a partial product by one word-size factor
at a time in the wider type. Equation~\eqref{eq:derivative} then gives the
derivative exactly. Operands zero and one produce zero directly.
Termination follows from strictly smaller recursive operands and finite
fallback trial division. An interrupt or unsuccessful arithmetic operation
leaves a block incomplete; it is never recorded as a negative result.

\subsection{Integer bounds and the square test}
For $x\ge2$, if $\Omega(x)=\sum e_p$, then
\begin{equation}\label{eq:bound}
\D(x)\le \frac{x\Omega(x)}2
\le \frac{x\lfloor\log_2x\rfloor}2.
\end{equation}
The first inequality follows from $p\ge2$ and the second from
$2^{\Omega(x)}\le x$. This is the elementary derivative bound underlying
the width analysis; its logarithmic version appears in~\cite[Theorem 9]{ufnarovski}.
The integer envelope $F(U)=\lfloor U\lfloor\log_2U\rfloor/2\rfloor$ is
nondecreasing for $U\ge2$. Starting from $U_0=9\cdot10^{14}$ and applying
$U_{j+1}=F(U_j)$ gives
\begin{align*}
U_1&=22050000000000000,\\
U_2&=595350000000000000,\\
U_3&=17562825000000000000<2^{64},\\
U_4&=553228987500000000000<2^{128}.
\end{align*}
The zero and one cases also satisfy these envelopes. Thus $x_0$ through
$x_3$ fit the word type, while $x_4$ uses the wider type. Even the conservative
bound $\sum U_j^2$ needs 138 bits: the square sum is therefore formed in GMP,
not in a fixed 128-bit accumulator. GMP's integer perfect-square predicate
and square root~\cite{gmp} decide~\eqref{eq:predicate}; the root is checked
on export to the wider type. There is no floating-point membership decision.
Indeed $r\le\sqrt{5}\,U_4<2^{71}$ in this envelope.

Additional program constraints support this domain on the measured 64-bit
platform. Blocks have a power-of-two width between $2^{14}$ and $2^{20}$,
thread counts are between one and eight, and segment counts are below $2^{40}$.
The prime table bound is at most $30000000$, fitting its 32-bit entries.
A block uses two 64-bit arrays and one byte array, at most 17 MiB altogether,
in addition to the prime table and worker state. Segment-width and endpoint
calculations use the wider type: a word-size block-group count times a
$2^{20}$ block is below $2^{84}$, and multiplying by a segment identifier
below $2^{40}$ remains below $2^{124}$, with the bounded endpoint additions
also safe. Allocation, I/O and elapsed-time errors cause failure rather than
certification. The cap is a sufficient numerical domain, not a maximum
mathematical domain, a runtime estimate, or a claim of feasible enumeration
to that endpoint. For comparison, the same sufficient envelope at $10^{15}$
already gives $U_3=19514250000000000000>2^{64}$.

\subsection{Parallel blocks and persistent coverage}
The residual sieve is followed by an OpenMP loop with dynamic scheduling.
Each input owns a separate result flag, and worker-local factorization and
GMP state avoid shared arithmetic mutation. A lock-free atomic stop flag
communicates interruption. All inputs must reach a complete positive or
negative state before the block is accepted. Positive flags are reevaluated
serially and emitted in increasing input order. A failed or interrupted
block does not advance the saved frontier.

Checkpoints persist only fully processed blocks. A segment is sealed only
when its frontier reaches its specified upper endpoint. Each saved state
contains its bounds, frontier and positive rows, together with a digest of
its serialized content; a separate seal binds the complete segment file.
Saved configuration binds the source, build recipe and executable identities.
Exclusive advisory locking prevents simultaneous cooperating writers, while
readers use shared locks. Atomic replacement uses an exclusive temporary
file, flushing, file synchronization, rename, and directory synchronization;
initialization also synchronizes the parent directory.

On restart, the program audits the stored records and resumes after the last
saved complete block. Work in an unsaved partial block is repeated. This
preserves coverage without promising resumability inside an input's
factorization. If a seal was committed before its aggregate tables, a valid
short aggregate prefix is classified as repairable and rebuilt from the
segment authority; a changed row is rejected. Within each run, sealed segment
endpoints and their adjacency establish coverage rather than an aggregate
count alone. The two campaign phases are separate runs; their cross-phase
join and combined cardinality are checked externally from the configurations
and completed records. These guarantees assume a cooperating local-filesystem environment.
They do not establish resistance to forged unkeyed digests, arbitrary hardware
faults, noncooperating writers, or every power-loss/filesystem behavior.

The native verification command audits identities, framing, bounds, seals,
ordering and aggregate consistency, and reevaluates the stored positive
witnesses with the same arithmetic backend. It does not process every
excluded input again. Accordingly, it corroborates persistent coverage and
positive consistency but is correlated with the enumeration kernel.

\subsection{Cost and reproducibility}
Sieving consecutive original inputs reduces repeated factorization of those
inputs. The derivative values themselves are not a consecutive interval.
The cost includes prime-table construction per segment, block sieving,
factorization of derived operands, exact square tests, and persistent-state
audits. In particular, deterministic trial division can dominate difficult
fallback cases. No global $O(N\log\log N)$ claim is made for the complete
predicate. Aggregate rebuilding also grows with the number of sealed records.

The source and build recipe specify C17 with optimization level three,
strict warnings, OpenMP, and linkage to GMP and OpenSSL. Independent positive
validators use Python integer arithmetic and exhaustive trial division,
followed by an integer square root and an exact equality test. The earlier
sixth-term validator must read the frozen one-hit state supplied to it;
the later two-hit state is a different protocol input. The seventh-term
validator imports that same trial-division routine, requires textual identity
of the recorded first hit with the frozen, previously validated sixth-term
input, and recomputes only the new seventh-term hit. These two validators are
independent of the C factorization backend, but share their own Python
arithmetic routine and the defining equation. The public deposit includes
the frozen inputs as well as the scripts; a script alone cannot replay
the historical protocol states.

\section{Validation and measured execution}\label{sec:validation}
The independent positive checks factored four operands for each new term,
computed the four derivatives, and tested the resulting sum with an exact
integer root. Each historical check requested and produced one distinct
new positive witness. For the sixth term the operand
factorizations are
\begin{align*}
x_0&=3\cdot43\cdot59\cdot144899,&x_1&=2^2\cdot102988867,\\
x_2&=2^4\cdot109\cdot199\cdot1187,&x_3&=2^4\cdot3\cdot23\cdot751909.
\end{align*}
For the seventh they are
\begin{align*}
x_0&=3\cdot7^7\cdot1597,&x_1&=7^8\cdot11\cdot83,\\
x_2&=2\cdot3\cdot7^7\cdot1327,&x_3&=7^7\cdot17\cdot859.
\end{align*}
The square sums are respectively
$6137488686385847025$ and $393395870932753190544$.
These checks verify membership, not the absence of smaller members.

Implementation checks retained in the project include the five baseline
singletons, each compared with the SymPy oracle, and independent
SymPy-based hit-set comparisons over
$[2,500]$, $[19150,19180]$ and $[129340,129370]$, respectively 499, 31 and
31 inputs. One-, four- and eight-thread comparisons on a low window and a
16384-input window ending at the technical cap tested deterministic output;
they are comparisons of the same program, not independent oracles for the
whole cap domain. Segment-grouping and interrupted-resume tests checked
the coverage contract, while altered rows, digests and frontiers were
required to fail. A forced-fallback harness exercised the factorization
fallback and wide-root paths.

These arithmetic and broad operational tests ran on the earlier build.
Their transfer to the campaign build is restricted to the source regions
that are byte-identical. The final change added an initialization path-length
guard; it did not change arithmetic, block processing or thread behavior.
Targeted final-build tests exercised that guard and adjacent I/O error paths.
AddressSanitizer and UndefinedBehaviorSanitizer checks on the earlier build
covered the first singleton and the small initial interval, with leak
sanitization unavailable. A separate Clang thread-sanitizer check exercised
a four-thread small interval with OpenMP instrumentation and a race-detection
control; it was not a whole-campaign or final-build thread-sanitizer run.
Final-build initialization sanitizer checks supplied additional bounded
evidence. These checks detect particular implementation
faults; none replaces the domain proof or independently enumerates billions
of negative cases.

During manuscript preparation, a bounded seed check verified the four small
factorizations in Theorem~\ref{thm:family} and their agreement with existing
rows 1, 2, 3 and 7. A bounded Python trial-division check also recomputed all
seven chains and their roots, independently of the C factorization backend;
it used the same Python factorization routine as the historical positive
validators. Another finite integer check verified the seven table
square identities, segment adjacency and counts, timing sum and width
envelopes. These are mathematical and bookkeeping checks, not extensions
of the official search.
The expanded exposition includes a bounded check of the eight prime
parameters $2,3,5,7,11,13,17,19$, the depth-two seed $128$, constant
fixed-point chains and three twin-prime examples. Direct trial-division
derivatives confirm the displayed positive chains. These checks corroborate
examples; the general statements depend on their proofs, and no new
sequence indices are certified.

The author executed the campaign on an AMD Ryzen 9 7940HS processor running
Linux Mint 22.3, with Linux kernel 6.17.0-42-generic, x86-64, glibc 2.39
and GCC 13.3.0.
The recorded runtime libraries include GMP, OpenSSL's crypto library and
GNU OpenMP. All seven externally timed segment-advance calls exited zero.
Their elapsed times sum to approximately 7984 seconds, or 2 hours 13 minutes
4 seconds, using eight threads. This total includes the work performed
inside those calls and excludes unrecorded initialization, separate audits
and manual activity. Dividing~\eqref{eq:count} by the recorded elapsed time
gives approximately $1.20\cdot10^6$ original inputs per second for this
campaign. This is an aggregate measurement, not a constant per-input cost
or an extrapolation to another decade or to the technical cap.

\section{Consequences and limits}
Relative to the accepted five-term baseline, the finite computation adds
two consecutive terms and excludes additional members through $10^{10}$.
Its negative evidence rests on complete execution of the exact C algorithm;
only positives and small windows have independent arithmetic comparisons.
The full earlier prefix interval was not rescanned.

The independent mathematical consequence is stronger than finite density
evidence: the explicit fixed-point family proves infinitude. It explains
four witnesses in the consolidated prefix while leaving three primitive
witnesses outside the family. Combining the proof and finite exclusion
gives~\eqref{eq:eighth}, but neither determines $a(8)$ nor rules out further
members between $10^{10}$ and the displayed upper bound. No conjecture about density,
successive indices or a complete classification is proposed.
The criterion in Section~\ref{sec:related} explains how the historical
mechanism yields families at other depths. The examples and source
comparison distinguish the mechanism from its particular specializations;
the literature search does not certify absolute priority.

\section{Data availability}
The public GitHub and Zenodo package~\cite{project} contains the source code,
build recipe, accepted baseline, indexed results, complete segment,
environment and timing records, independent positive-check scripts with
frozen inputs, bounded seed checks and validation summaries. The manuscript
is supplied as monolithic LaTeX source and a matching PDF; its witness table
is embedded here. Both providers' initial release payloads were downloaded,
checked against the full inventory and independently passed the bounded
validation command. This does not independently rescan negative intervals.
The designated DOI identifies the release series; exact deposits have
separate DOIs. Primary-source locators are included; third-party full texts,
private research notes and external reviews are excluded.

\section{AI use disclosure}
Large language models were used as research and drafting assistants during
the A249106 workflow, including OpenAI Codex and external LLM-based review
systems from the ChatGPT, Claude, DeepSeek, Gemini, Grok, Qwen, MiniMax,
Meta and Z.ai families. They are named at family level where exact
model-version metadata were not consistently available in the local artifacts.
\paragraph{Mathematical analysis.}
OpenAI Codex assisted in reconstructing the definition, checking primary
sources, analyzing the seed scaling identity, and organizing the proof and
finite evidence. AI-generated preliminary triage reports and a stored
NotebookLM synthesis informed earlier selection and strategy research;
they were not treated as proofs or as substitutes for primary sources.
One preliminary triage had noted the shared factor $1597$ in the first three
terms as an unproved pattern; the fixed-point family and its proof followed
the finite campaign.

\paragraph{Algorithm and code.}
OpenAI Codex assisted with implementation and local validation. The retained
code-review reports include OpenAI, Anthropic Claude, Z.ai GLM-5.3, Google
Gemini, xAI Grok, DeepSeek, Meta, MiniMax and Alibaba Qwen tools. Their roles
were code and artifact review, with execution claims distinguished from
static reasoning. Model versions are not inferred from filenames or behavior;
the attribution of GLM-5.3 is to Z.ai. These reports concern code and do not
constitute external review of this manuscript.

\paragraph{Manuscript.}
OpenAI Codex assisted in drafting the English text and in sequential internal
mathematical, numerical, source and exposition checks. These passes share the
drafting context and are not independent external reviews. The first full
external manuscript round supplied nine reports from the ChatGPT, Claude,
DeepSeek, Gemini, GLM, Grok, Meta, MiniMax and Qwen families. Their observations
were checked against the source artifacts before incorporation; reported
remote executions were not treated as locally reproduced evidence.
A focused second round supplied reports from Claude and ChatGPT, leading
to final clarifications of positive-validation provenance and citation wording.
Family overlap with code review, and the OpenAI drafting/review relationship,
limit independence.
The author subsequently requested a fuller account of Theorem~\ref{thm:family}.
OpenAI Codex assisted with the primary-source comparison, the general seed
criterion and the A230543 example, followed by bounded arithmetic checks
and sequential internal review. A focused third review of the expanded
text, from the Claude family, found no scientific issue and suggested
clarifications of attribution, scope and exposition. These were checked
against the retained sources before incorporation; family overlap with
earlier rounds limits independence. The final claims
are based on source artifacts, recorded computations, independent positive
validators and primary literature, not conversational assertions.
Carlo Corti retains full scientific responsibility for the mathematical
statements, implementation, numerical results, validation boundaries and
conclusions.

\clearpage
\begin{thebibliography}{99}
\bibitem{sequence}
P. P. Lava, OEIS A249106, Numbers that form a Pythagorean 6-tuple with their
first four arithmetic derivatives (2014), with extensions by R. Chandler
(2016, 2017). \url{https://oeis.org/A249106}.
Consulted via the archived entry snapshot matching the public entry.

\bibitem{derivative}
N. J. A. Sloane and R. K. Guy, OEIS A003415 (arithmetic derivative of $n$).
\url{https://oeis.org/A003415}. Consulted October 3, 2026.

\bibitem{auxiliary}
A. Auel, OEIS A051674, $p^p$ for prime $p$;
R. Zumkeller, OEIS A129283, $n+\D(n)$ (April 7, 2007).
\url{https://oeis.org/A051674}, \url{https://oeis.org/A129283}.
Consulted October 3, 2026.

\bibitem{related}
P. P. Lava, OEIS entries A210503 (January 25, 2013),
A230543 (October 25, 2013), and A249105, A249107 and A249110
(October 21, 2014).
\url{https://oeis.org/A210503}, \url{https://oeis.org/A230543},
\url{https://oeis.org/A249105}, \url{https://oeis.org/A249107},
\url{https://oeis.org/A249110}. Consulted October 3, 2026.

\bibitem{ufnarovski}
V. Ufnarovski and B. \AA hlander, How to differentiate a number,
\emph{Journal of Integer Sequences} \textbf{6} (2003), Article 03.3.4.
\url{https://cs.uwaterloo.ca/journals/JIS/VOL6/Ufnarovski/ufnarovski.pdf}.

\bibitem{kovic}
J. Kovi\v{c}, The arithmetic derivative and antiderivative,
\emph{Journal of Integer Sequences} \textbf{15} (2012), Article 12.3.8.
\url{https://cs.uwaterloo.ca/journals/JIS/VOL15/Kovic/kovic4.html}.
Original full-text source consulted October 3, 2026.

\bibitem{frisch}
S. Frisch and L. Vaserstein, Polynomial parametrization of Pythagorean
quadruples, quintuples and sextuples, \emph{Journal of Pure and Applied
Algebra} \textbf{216} (2012), 184--191.
\url{https://doi.org/10.1016/j.jpaa.2011.06.002}.
Text consulted: arXiv:1106.3493v2 (June 28, 2011).
Preprint DOI: \url{https://doi.org/10.48550/arXiv.1106.3493}.
Theorem numbering refers to that preprint version.

\bibitem{forisek}
M. Fori\v{s}ek and J. Jan\v{c}ina, Fast primality testing for integers that fit
into a machine word, in \emph{Student Research Forum Papers and Posters at
SOFSEM 2015}, CEUR Workshop Proceedings \textbf{1326} (2015), 20--30.
\url{https://ceur-ws.org/Vol-1326/020-Forisek.pdf}.

\bibitem{pollard}
J. M. Pollard, A Monte Carlo method for factorization,
\emph{BIT} \textbf{15} (1975), 331--334.
\url{https://doi.org/10.1007/BF01933667}.

\bibitem{gmp}
GNU MP project, GNU MP 6.3.0 manual, Section 5.8, Root extraction functions.
\url{https://gmplib.org/manual/Integer-Roots}. Consulted October 2, 2026.

\bibitem{project}
C. Corti, \emph{An Exact Extension and an Infinite Family for OEIS A249106}
(2026), GitHub repository and Zenodo archive.
\url{https://github.com/carcorti/A249106};
\url{https://doi.org/10.5281/zenodo.23155533}.
The Zenodo DOI identifies the release series; exact deposited versions have
separate version-specific DOIs.
\end{thebibliography}
\end{document}

